\documentclass[a4paper,12pt]{article}

\usepackage{fontspec}
\usepackage{polyglossia}
\setdefaultlanguage{russian}
\setotherlanguage{english}
\setmainfont{Times New Roman}
\setsansfont{Arial}
\setmonofont{Menlo}
\newfontfamily\cyrillicfont{Times New Roman}
\newfontfamily\cyrillicfontsf{Arial}
\newfontfamily\cyrillicfonttt{Menlo}
\usepackage{amsmath,amsfonts,amssymb}
\usepackage{graphicx}
\usepackage{color}
\usepackage{hyperref}
\usepackage{geometry}
\usepackage{indentfirst}
\usepackage{longtable}
\usepackage{booktabs}
\usepackage{float}
\usepackage{listings}
\usepackage{xcolor}
\geometry{left=3cm,right=1.5cm,top=2cm,bottom=2cm}

\lstset{
    basicstyle=\ttfamily\small,
    breaklines=true,
    frame=single,
    backgroundcolor=\color{gray!10},
    keywordstyle=\color{blue},
    commentstyle=\color{green!60!black},
    stringstyle=\color{red!70!black},
    showstringspaces=false,
    tabsize=4
}

\begin{document}

\begin{titlepage}
    \centering
    {\fontsize{14pt}{16pt}\selectfont МОСКОВСКИЙ АВИАЦИОННЫЙ ИНСТИТУТ\\(НАЦИОНАЛЬНЫЙ ИССЛЕДОВАТЕЛЬСКИЙ УНИВЕРСИТЕТ)}\\[0.5cm]
    {\fontsize{12pt}{14pt}\selectfont Институт №8 «Компьютерные науки и прикладная математика»}\\[4cm]
    
    {\fontsize{16pt}{18pt}\selectfont \textbf{Лабораторные работы}}\\
    {\fontsize{14pt}{16pt}\selectfont \textbf{по курсу «Информационный поиск»}}\\[5cm]
    
    \vfill
    \begin{flushright}
        \begin{minipage}{0.51\textwidth}
            Выполнил: Горюнов Д.В. \\
            Группа: М8О-408Б-22\\
            Преподаватель: Кухтичев Антон Алексеевич
        \end{minipage}
    \end{flushright}
    
    \vfill
    {\large Москва, 2026}
\end{titlepage}

\tableofcontents
\newpage

\section{Лабораторная работа №1. Добыча корпуса документов}

\subsection{Задание}
Необходимо проанализировать корпус документов, который будет использован при выполнении 
остальных лабораторных работ: скачать документы, ознакомиться с их характеристиками, 
выделить текст, найти существующие поисковики для данного набора документов.

\subsection{Метод решения}

\textbf{Тема корпуса:} Автомобильные новости и обзоры на английском языке (Automotive news and reviews).

\textbf{Источники данных:}
\begin{enumerate}
    \item \textbf{The Drive} (\url{https://www.thedrive.com}) --- крупнейший автомобильный 
    медиа-ресурс с тысячами новостей, обзоров и тест-драйвов (32\,870 статей).
    \item \textbf{TopSpeed} (\url{https://www.topspeed.com}) --- автомобильный портал 
    с обзорами и сравнениями автомобилей (2\,237 статей).
    \item \textbf{CarBuzz} (\url{https://carbuzz.com}) --- новостной автомобильный портал (818 статей).
    \item \textbf{CarScoops} (\url{https://www.carscoops.com}) --- автомобильные новости (31 статья).
\end{enumerate}

Документы представляют собой HTML-страницы с автомобильными статьями. Каждый документ 
содержит:
\begin{itemize}
    \item Заголовок статьи (тег \texttt{<h1>});
    \item Основной текст статьи внутри тега \texttt{<article>};
    \item Мета-информацию: дата публикации, автор, категория;
    \item Навигационные элементы, рекламу, скрипты (удаляются при извлечении).
\end{itemize}

Для извлечения текста из HTML используется скрипт на Python с библиотекой BeautifulSoup. 
Скрипт удаляет теги \texttt{script}, \texttt{style}, \texttt{nav}, \texttt{footer}, 
\texttt{header}, \texttt{aside}, после чего извлекает чистый текст.

\textbf{Существующие поисковики для данного корпуса:}
\begin{itemize}
    \item Встроенный поиск на сайтах thedrive.com и topspeed.com;
    \item Google с ограничением: \texttt{site:thedrive.com electric vehicle review}.
\end{itemize}

\subsection{Результаты}

\begin{figure}[H]
    \centering
    \includegraphics[width=0.75\textwidth]{screenshots/corpus_stats.png}
    \caption{Статистика корпуса документов}
    \label{fig:corpus_stats}
\end{figure}

Корпус содержит 35\,889 документов из четырёх источников. Общий объём сырого HTML --- 
4\,979~МБ. Общий объём извлечённого текста составляет 140.7~МБ. Средний размер 
извлечённого документа --- 4\,109 байт.

\begin{figure}[H]
    \centering
    \includegraphics[width=0.95\textwidth]{screenshots/raw_vs_clean.png}
    \caption{Сравнение сырого HTML и извлечённого текста}
    \label{fig:raw_vs_clean}
\end{figure}

На рис.~\ref{fig:raw_vs_clean} показано сравнение сырого HTML-документа с извлечённым 
чистым текстом. Извлечение существенно уменьшает объём данных, сохраняя при этом весь 
информативный контент.

\subsection{Выводы}
Корпус из 35\,889 реальных автомобильных статей на английском языке успешно собран из 
четырёх независимых источников. Средний объём извлечённого текста 4\,109 байт, 
что соответствует нескольким сотням слов --- достаточно для качественной 
токенизации и индексирования.

\newpage

\section{Лабораторная работа №2. Поисковый робот}

\subsection{Задание}
Написать поисковый робот на любом языке программирования. Робот должен принимать YAML-конфиг, 
сохранять документы в базу данных с полями url, html, source, дата обкачки. Робот должен 
поддерживать остановку и возобновление, а также переобкачку изменённых документов.

\subsection{Метод решения}

Поисковый робот реализован на Python с использованием асинхронного подхода (\textbf{aiohttp})
для параллельной загрузки. Также реализован вариант на \textbf{Scrapy}.
В качестве базы данных используется \textbf{SQLite}.

\textbf{Конфигурация} задаётся в файле \texttt{config.yaml}:
\begin{lstlisting}[language=Python]
db:
  path: "../data/corpus.db"
logic:
  delay: 0.5
  user_agent: "IRLabBot/1.0"
  max_pages: 20000
  concurrent_requests: 4
\end{lstlisting}

\textbf{Схема базы данных SQLite:}
\begin{lstlisting}[language=SQL]
CREATE TABLE documents (
    id INTEGER PRIMARY KEY AUTOINCREMENT,
    url TEXT UNIQUE NOT NULL,
    html TEXT NOT NULL,
    title TEXT,
    source TEXT NOT NULL,
    crawled_at INTEGER NOT NULL,
    content_hash TEXT NOT NULL
);
\end{lstlisting}

\textbf{Архитектура краулера:}
\begin{itemize}
    \item Фаза 1: обход sitemap.xml каждого сайта, сбор URL статей;
    \item Фаза 2: параллельная загрузка (60--80 воркеров) через aiohttp;
    \item Сохранение документов в SQLite с вычислением MD5-хеша контента;
    \item \textbf{Возобновление:} при запуске проверяются уже скачанные URL (UNIQUE constraint), 
    дубликаты пропускаются;
    \item \textbf{Переобкачка:} если хеш контента изменился, документ обновляется в базе.
\end{itemize}

\subsection{Результаты}

\begin{figure}[H]
    \centering
    \includegraphics[width=0.95\textwidth]{screenshots/crawler_run.png}
    \caption{Лог работы поискового робота}
    \label{fig:crawler_run}
\end{figure}

Робот использует асинхронную загрузку с 60--80 параллельными воркерами. 
Скорость обкачки --- до 20 страниц/с. 
Общее время обкачки 35\,956 документов из 4 источников --- около 75 минут.
Размер базы данных --- 5~ГБ (сырой HTML).

\subsection{Выводы}
Реализованный поисковый робот с асинхронной загрузкой обеспечивает скоростную обкачку 
автомобильных новостных сайтов. Использование SQLite упрощает развёртывание --- не 
требуется отдельный сервер БД. Механизм хеширования контента позволяет эффективно 
определять изменённые документы при повторном обходе.

\newpage

\section{Лабораторная работа №3. Токенизация}

\subsection{Задание}
Реализовать процесс разбиения текстов документов на токены. Описать правила токенизации, 
указать достоинства и недостатки выбранного метода.

\subsection{Метод решения}

Токенизатор реализован на C++ (STL разрешена для данной лабораторной). 

\textbf{Правила токенизации:}
\begin{enumerate}
    \item Текст разбивается по не-буквенным символам (пробелы, знаки пунктуации, цифры, 
    спецсимволы);
    \item Все символы приводятся к нижнему регистру (\texttt{tolower});
    \item Дефисы и апострофы: апостроф сохраняется внутри слова (\texttt{don't} $\to$ 
    \texttt{don't}), дефис разделяет токены;
    \item Токены короче 2 символов отбрасываются;
    \item Числовые токены игнорируются.
\end{enumerate}

\textbf{Достоинства метода:}
\begin{itemize}
    \item Простота реализации и высокая скорость работы;
    \item Корректная обработка основных случаев для английского текста;
    \item Сохранение апострофов позволяет различать формы (\texttt{it's}, \texttt{don't}).
\end{itemize}

\subsection{Результаты}

\begin{figure}[H]
    \centering
    \includegraphics[width=0.85\textwidth]{screenshots/tokenization.png}
    \caption{Пример работы токенизатора}
    \label{fig:tokenization}
\end{figure}

\begin{table}[H]
\centering
\begin{tabular}{lc}
\toprule
\textbf{Метрика} & \textbf{Значение} \\
\midrule
Количество токенов (весь корпус) & 23\,568\,294 \\
Средняя длина токена & 4.79 символов \\
Среднее количество токенов на документ & 656 \\
Скорость токенизации & 20\,598.1 КБ/с \\
\bottomrule
\end{tabular}
\caption{Статистика токенизации}
\label{tab:tokenization_stats}
\end{table}

Скорость токенизации составляет 20\,598~КБ/с, что является высоким показателем. 
Обработка всего корпуса (140.7~МБ текста) занимает около 7 секунд.

\subsection{Выводы}
Реализованный токенизатор корректно обрабатывает английский текст автомобильной тематики. 
Средняя длина токена 4.79~символов соответствует ожидаемым значениям для английского языка. 
Скорость токенизации (более 20~МБ/с) достаточна для обработки корпусов масштаба 
сотен тысяч документов.

\newpage

\section{Лабораторная работа №4. Закон Ципфа}

\subsection{Задание}
Построить график распределения терминов по частотностям в логарифмической шкале, наложить 
на него закон Ципфа. Объяснить причины расхождения.

\subsection{Метод решения}

Для построения графика используется следующий подход:
\begin{enumerate}
    \item С помощью C++ утилиты \texttt{search\_engine freq} вычисляются частоты всех термов 
    (после стемминга) в корпусе;
    \item Python-скрипт \texttt{zipf.py} строит график в лог-лог шкале с помощью matplotlib;
    \item На график накладывается идеальная кривая Ципфа: $f(r) = C / r$, где $C$ --- 
    частота самого распространённого терма, $r$ --- ранг.
\end{enumerate}

\subsection{Результаты}

\begin{figure}[H]
    \centering
    \includegraphics[width=0.85\textwidth]{screenshots/zipf.png}
    \caption{Распределение терминов по частотностям и закон Ципфа}
    \label{fig:zipf}
\end{figure}

Всего в корпусе выделено 80\,453 уникальных терма (после стемминга). Наиболее частый 
терм --- <<the>> (1\,433\,009 вхождений). График показывает, что реальное распределение 
близко к закону Ципфа, но не совпадает с ним точно.

\textbf{Причины расхождения:}
\begin{itemize}
    \item \textbf{В области высоких рангов} (редкие слова): частота убывает медленнее, чем 
    предсказывает закон Ципфа. Это связано с тематической специализацией корпуса --- 
    автомобильная лексика содержит много технических терминов с умеренной частотой;
    \item \textbf{В области низких рангов} (частые слова): наиболее частые слова имеют 
    частоту несколько выше предсказанной, что объясняется узкой тематикой корпуса;
    \item Стемминг объединяет словоформы, изменяя распределение частот по сравнению с 
    исходными токенами.
\end{itemize}

\subsection{Выводы}
Распределение терминов по частотностям в нашем корпусе в целом следует закону Ципфа, 
что подтверждает его применимость для тематических корпусов. Отклонения объясняются 
спецификой автомобильной тематики и эффектом стемминга.

\newpage

\section{Лабораторная работа №5. Стемминг}

\subsection{Задание}
Добавить в поисковую систему стемминг.

\subsection{Метод решения}

Реализован алгоритм Портера (Porter Stemmer) для английского языка на C++ без использования 
STL. Алгоритм состоит из 5 шагов последовательного удаления и замены суффиксов.

\textbf{Шаги алгоритма:}
\begin{enumerate}
    \item \textbf{Шаг 1a}: обработка множественного числа (\texttt{-sses} $\to$ \texttt{-ss}, 
    \texttt{-ies} $\to$ \texttt{-i}, \texttt{-s} $\to$ удаление);
    \item \textbf{Шаг 1b}: обработка прошедшего времени и герундия (\texttt{-eed}, \texttt{-ed}, 
    \texttt{-ing});
    \item \textbf{Шаг 1c}: замена конечного \texttt{y} на \texttt{i};
    \item \textbf{Шаг 2}: замена суффиксов (\texttt{-ational} $\to$ \texttt{-ate}, 
    \texttt{-izer} $\to$ \texttt{-ize} и др.) при условии $m > 0$;
    \item \textbf{Шаг 3}: удаление суффиксов (\texttt{-icate} $\to$ \texttt{-ic}, 
    \texttt{-ful} $\to$ удаление и др.);
    \item \textbf{Шаг 4}: удаление суффиксов при $m > 1$ (\texttt{-al}, \texttt{-ance}, 
    \texttt{-ence}, \texttt{-ion} и др.);
    \item \textbf{Шаг 5}: финальная обработка (удаление конечного \texttt{e}, дублированного 
    \texttt{l}).
\end{enumerate}

Стеммер работает с собственной реализацией строки (\texttt{MyString}) без использования STL.

\subsection{Результаты}

\begin{figure}[H]
    \centering
    \includegraphics[width=0.65\textwidth]{screenshots/stemming.png}
    \caption{Примеры работы Porter Stemmer}
    \label{fig:stemming}
\end{figure}

Примеры на рис.~\ref{fig:stemming} демонстрируют корректную работу стеммера на 
автомобильной лексике: \texttt{running}~$\to$~\texttt{run}, 
\texttt{acceleration}~$\to$~\texttt{acceler}, \texttt{turbocharged}~$\to$~\texttt{turbocharg}.

Стемминг уменьшает количество уникальных терминов, что повышает полноту поиска (recall). 
Например, запрос \texttt{drive} теперь находит документы, содержащие \texttt{driving}, 
\texttt{driven}, \texttt{driver}.

\subsection{Выводы}
Алгоритм Портера корректно обрабатывает основные морфологические формы английских слов. 
Интеграция стеммера в поисковую систему повышает полноту поиска за счёт приведения различных 
словоформ к единой основе.

\newpage

\section{Лабораторная работа №6. Булев индекс}

\subsection{Задание}
Построить поисковый индекс, пригодный для булева поиска: самостоятельно разработанный 
бинарный формат, обратный и прямой индексы.

\subsection{Метод решения}

\subsubsection{Бинарный формат инвертированного индекса}

Файл \texttt{index.bin} имеет следующую структуру:

\begin{lstlisting}
[Magic 4B: "BIDX"] [Version 2B] [NumTerms 4B] [NumDocs 4B]
For each term:
  [term_len 2B] [term_chars ...] [doc_freq 4B]
  [posting list: doc_freq * 4B doc_id values]
\end{lstlisting}

\begin{figure}[H]
    \centering
    \includegraphics[width=0.9\textwidth]{screenshots/index_format.png}
    \caption{Структура бинарного формата индекса}
    \label{fig:index_format}
\end{figure}

\subsubsection{Прямой индекс}

Файл \texttt{forward.bin}:
\begin{lstlisting}
[Magic 4B: "FIDX"] [NumDocs 4B]
For each document:
  [url_len 2B] [url_chars ...] [title_len 2B] [title_chars ...]
\end{lstlisting}

\subsubsection{Ключевые структуры данных}

Все структуры данных реализованы вручную на C++ без использования STL:
\begin{itemize}
    \item \textbf{MyString} --- строка (\texttt{char*}, \texttt{len}, \texttt{cap}) с 
    операциями push, append, tolower, compare;
    \item \textbf{DynArray / UInt32Array} --- динамический массив с автоматическим 
    расширением (удвоение ёмкости);
    \item \textbf{HashMap} --- хеш-таблица с функцией \textbf{FNV-1a} и открытой адресацией 
    (линейное пробирование), автоматическое расширение при заполнении 70\%.
\end{itemize}

\textbf{Хеш-функция FNV-1a:}
\begin{lstlisting}[language=C++]
uint32_t fnv1a_hash(const char* data, size_t len) {
    uint32_t hash = 2166136261u;
    for (size_t i = 0; i < len; i++) {
        hash ^= (uint8_t)data[i];
        hash *= 16777619u;
    }
    return hash;
}
\end{lstlisting}

\textbf{Метод сортировки:} QuickSort (быстрая сортировка) для упорядочивания posting lists. 
Средняя сложность $O(n \log n)$, хорошо работает на случайных данных.

\subsection{Результаты}

\begin{figure}[H]
    \centering
    \includegraphics[width=0.9\textwidth]{screenshots/index_build.png}
    \caption{Процесс построения индекса}
    \label{fig:index_build}
\end{figure}

\begin{table}[H]
\centering
\begin{tabular}{lc}
\toprule
\textbf{Метрика} & \textbf{Значение} \\
\midrule
Количество документов & 35\,889 \\
Количество уникальных термов & 80\,453 \\
Размер инвертированного индекса & 42 МБ \\
Размер прямого индекса & 911 КБ \\
Время построения & 6.99 с \\
Скорость индексации & 5\,132.7 док/с \\
Скорость индексации & 20\,598.1 КБ/с \\
\bottomrule
\end{tabular}
\caption{Статистика индексации}
\label{tab:index_stats}
\end{table}

Средняя длина терма после стемминга составляет примерно 4.5 символа, что меньше средней 
длины токена (4.79). Это объясняется тем, что стемминг удаляет суффиксы, сокращая 
длину слов.

\subsection{Выводы}
Бинарный формат индекса обеспечивает компактное хранение и быстрое чтение. Хеш-таблица 
с FNV-1a и открытой адресацией обеспечивает $O(1)$ доступ к posting list для каждого терма. 
Скорость индексации позволяет обрабатывать корпус за несколько секунд. При увеличении 
объёма данных в 10 раз основным ограничением станет объём оперативной памяти для хеш-таблицы; 
в этом случае потребуется переход к блочной индексации (BSBI или SPIMI).

\newpage

\section{Лабораторная работа №7. Булев поиск}

\subsection{Задание}
Реализовать ввод поисковых запросов и их выполнение над индексом. Синтаксис запросов: 
пробел или \texttt{\&\&} для И, \texttt{||} для ИЛИ, \texttt{!} для НЕТ, поддержка скобок. 
Реализовать веб-сервис с формой ввода и страницей выдачи.

\subsection{Метод решения}

\subsubsection{Парсер запросов}

Парсер реализован методом рекурсивного спуска на C++ без STL. Грамматика:

\begin{lstlisting}
expr     -> or_expr
or_expr  -> and_expr ("||" and_expr)*
and_expr -> not_expr (("&&" | implicit_and) not_expr)*
not_expr -> "!" not_expr | atom
atom     -> "(" expr ")" | WORD
\end{lstlisting}

Парсер устойчив к переменному числу пробелов. Пробел между словами трактуется как 
операция И (implicit AND).

\subsubsection{Выполнение запросов}

Каждый терм ищется в хеш-таблице инвертированного индекса. Posting lists (отсортированные 
массивы doc\_id) объединяются операциями:
\begin{itemize}
    \item \textbf{AND (пересечение)}: два указателя проходят по двум массивам, $O(n + m)$;
    \item \textbf{OR (объединение)}: слияние двух отсортированных массивов, $O(n + m)$;
    \item \textbf{NOT}: дополнение множества до полного набора документов, $O(N)$ где $N$ --- 
    число документов.
\end{itemize}

\subsubsection{Веб-интерфейс}

Веб-сервис реализован на Python с использованием встроенного \texttt{http.server}:
\begin{itemize}
    \item \textbf{GET /} --- начальная страница с формой ввода запроса;
    \item \textbf{GET /search?q=...} --- страница результатов с 50 результатами на странице, 
    содержащая заголовки документов и ссылки;
    \item Веб-сервер вызывает C++ CLI через \texttt{subprocess}.
\end{itemize}

\subsection{Результаты}

\begin{figure}[H]
    \centering
    \includegraphics[width=0.9\textwidth]{screenshots/search_results.png}
    \caption{Результаты булева поиска}
    \label{fig:search_results}
\end{figure}

\begin{table}[H]
\centering
\begin{tabular}{lcc}
\toprule
\textbf{Запрос} & \textbf{Результатов} & \textbf{Время, мс} \\
\midrule
\texttt{toyota camry} & 401 & 0.08 \\
\texttt{electric vehicle} & 5\,978 & 0.47 \\
\texttt{tesla || rivian} & 4\,071 & 0.11 \\
\texttt{bmw !audi} & 3\,209 & 0.47 \\
\texttt{(ford || chevrolet) truck} & 3\,573 & 0.46 \\
\texttt{hybrid suv review} & 1\,058 & 0.52 \\
\bottomrule
\end{tabular}
\caption{Скорость выполнения поисковых запросов}
\label{tab:search_speed}
\end{table}

\begin{figure}[H]
    \centering
    \includegraphics[width=0.95\textwidth]{screenshots/web_interface.png}
    \caption{Веб-интерфейс поисковой системы}
    \label{fig:web_interface}
\end{figure}

Наиболее длительный запрос --- \texttt{hybrid suv review} (0.52~мс), так как он 
включает пересечение трёх posting lists. Простые запросы выполняются за 0.08--0.11~мс.
Операция NOT (\texttt{bmw !audi}) требует перебора всех 35\,889 документов --- 0.47~мс.

\textbf{Тестирование корректности} проводилось путём сравнения результатов булева поиска 
с ручной проверкой: для каждого тестового запроса проверялось, что возвращённые документы 
действительно содержат (или не содержат) указанные термины.

\subsection{Выводы}
Реализованная система булева поиска обеспечивает выполнение запросов за доли миллисекунды 
благодаря использованию отсортированных posting lists и эффективных операций пересечения 
и объединения. Веб-интерфейс обеспечивает удобное взаимодействие с поисковой системой через 
браузер. CLI-утилита позволяет выполнять пакетные запросы из файла.

\end{document}
